\section{Vision and approach}

\subsection{Vision}

From autonomous drone swarms navigating disaster zones to crowds flowing through transport hubs, understanding how groups coordinate and compete is fundamental to critical societal challenges. Current mathematical methods for analysing these systems have a practical limitation: they examine behaviour at a single spatial scale, missing the hierarchical organisation, individual decisions triggering tactical responses that shape system-wide patterns, characteristic of most real-world competitive multi-agent systems. We have developed and validated a multi-scale persistent homology framework addressing this limitation (Brown et al., 2026, ArXiv preprint). This research scales the validated framework from proof of concept to full season, developing persistence-landscape methods for temporal dynamics, using professional football as a testbed in which high-quality data enables rigorous method development.

\paragraph{Excellence and importance.}
Topological Data Analysis (TDA) has proven effective for characterising the shape of complex data, but existing applications examine systems at single scales, in non-competitive scenarios, or with static data. A standard Vietoris--Rips filtration on a multi-agent point cloud conflates topological features from different organisational levels into a single persistence diagram. Our framework addresses this through two methodological contributions: domain-informed hierarchical clustering to decompose point clouds by organisational level, and adaptive filtration ensuring consistent $H_1$ (loop) detection across all scales, a practical necessity since a fixed Vietoris--Rips threshold appropriate at one level produces null results at another. These contributions sit within computational topology: practical methods for multi-scale persistent homology on hierarchical, competitive, time-evolving point clouds.

Validation across 11 professional football matches (1{,}650 analysis frames, 1 optical tracking + 10 broadcast tracking) confirms three robust $H_0$ analysis regimes, individual (3~m), tactical (12~m), and team (30~m), with stability scores exceeding 0.88, and two $H_1$ regimes capturing complementary structural information. Individual-scale $H_1$ reveals frequent, transient loops ($97.2\% \pm 1.6\%$ presence rate across matches), whilst tactical-scale $H_1$ reveals rarer but geometrically distinct loops ($20.1\% \pm 7.6\%$ presence rate). Scale complementarity is quantified: Spearman $\rho = 0.25$ between scales, confirming that neither subsumes the other. Real event correlation across 10 matches (104{,}722 event--topology pairs) demonstrates that topological features respond coherently to match dynamics, disruption events decrease persistence, organisational events increase it. Temporal persistence dynamics are match-specific rather than universal, unsurprising to practitioners given the diversity of tactical systems, match states, and team strategies, but the framework provides the first tools to quantify this variation systematically.

\paragraph{Research objectives.}
Building on this validated foundation, three research objectives are pursued over 12 months, with a fourth strategic output: (1)~Scale the analysis pipeline to a full Championship season ($\sim$540 matches), stress-testing the framework and establishing population-level topological statistics; (2)~Characterise pattern formation by developing baseline topological fingerprints for different tactical systems, quantifying formation times and distinguishing between tactical configurations in persistence space; (3)~Develop persistence landscape methods (Bubenik, 2015) for temporal dynamics, classifying stability regimes and characterising transitions within and across matches; (4)~Synthesise full-season results into a Standard Grant application extending the framework to broader competitive systems.

\paragraph{Advancing current understanding.}
The 11-match validation establishes that multi-scale topological structure is a robust feature of competitive spatial systems, not an artefact of any single match. The framework now addresses open questions in applied computational topology: how does one perform meaningful scale-specific persistent homology when the system has known hierarchical organisation? Our answer, clustering-then-filtration with adaptive thresholds, is general, reproducible, and validated on real data. This grant extends this in three directions: (i)~full-season analysis establishes whether topological signatures are stable population features or exhibit meaningful distributional variation; (ii)~tactical fingerprinting tests whether persistence diagrams can discriminate between tactical systems, connecting TDA to practical classification problems; (iii)~persistence landscape dynamics develops functional representations of temporal evolution, creating the mathematical apparatus needed for stability regime characterisation. Publications target the \emph{Journal of Applied and Computational Topology} and \emph{SIAM Journal on Applied Mathematics}.

\paragraph{Timeliness.}
Computational advances in TDA libraries (Ripser, GUDHI) now enable real-time computation on large point clouds, making full-season multi-scale analysis feasible. Theoretical momentum is building around persistent homology for dynamical systems, with growing interest in temporal methods and stability theory. High-quality spatial tracking data from professional football provides an ideal testbed, offering high-frequency measurements (10--25~Hz) on a well-structured competitive system. Industry demand is demonstrated through established partnerships with championship clubs. The framework aligns with UKRI priorities for explainable AI, since TDA features are interpretable by construction: each topological feature has a direct geometric realisation, in contrast to black-box machine learning approaches.

\paragraph{Impact.}
Direct beneficiaries include the mathematical sciences community (validated methods for multi-scale TDA in competitive systems), sports science researchers (quantitative tools for spatial analysis), and industry partners (championship clubs). Preliminary cross-domain application to armed conflict event data (Brown et al., in preparation) suggests the three-scale structure transfers across domains with appropriate scale adjustment, defining boundary conditions for when the framework produces meaningful results. The Small Grant provides the full-season evidence base needed to underpin a cross-domain Standard Grant. Indirect beneficiaries include autonomous systems developers (multi-agent coordination), emergency services (crowd management), and conservation organisations (ecosystem monitoring).

\paragraph{Pathway to larger research programme.}
This Small Grant serves as a stepping stone to a Standard Grant by generating full-season validation, peer-reviewed publications, and validated industry partnerships. The ArXiv preprint (Brown et al., 2026) establishes the methodology; this grant produces the scaling evidence. A Standard Grant submission is planned within 6--8 months of project completion.

\subsection{Approach}

\paragraph{Design for effective achievement of objectives.}
Three research objectives with clear success criteria, plus a strategic fourth output:

\textit{Objective 1, Full-Season Analysis (Months 1--7).} Scale the validated pipeline to the full Championship season ($\sim$540 matches). PI establishes the concurrent processing pipeline on Supercomputing Wales (Months 1--2) and analyses an initial cohort of 15--20 matches (Months 1--3). RA executes full-season analysis of remaining matches (Months 3--7). Success criteria: ${>}500$ matches processed; population-level $H_0$ and $H_1$ distributions established; cross-match variance characterised.

\textit{Objective 2, Tactical Fingerprinting (Months 4--8).} Develop baseline topological fingerprints for different tactical systems. Using formation metadata from StatsBomb events and broadcast annotation, classify matches by tactical system and compare their persistence diagram distributions. Success criteria: topological signatures for $\geq 3$ tactical configurations (e.g.\ 4-4-2, 4-3-3, 3-5-2); statistical tests for inter-system differences ($p < 0.05$ with FDR correction).

\textit{Objective 3, Persistence Landscape Dynamics (Months 6--10).} Implement persistence landscape representations (Bubenik, 2015) for temporal evolution analysis. Classify stability regimes within matches, characterise transitions, and test whether landscape dynamics correlate with tactical changes. Success criteria: functional persistence landscape computation at both analysis scales; identified stability regimes; transition characterisation.

\textit{Strategic Output, Standard Grant Preparation (Months 9--12).} Synthesise full-season results, persistence landscape theory, and tactical fingerprinting evidence into a Standard Grant proposal. Identify the key mathematical questions emerging from full-season data that require a larger programme to address.

\paragraph{Feasibility and risk management.}
The project is highly feasible. The framework has been validated across 11 matches with computational efficiency (${<}2$\,s per frame on Supercomputing Wales), and the full pipeline, from data loading through persistence computation to statistical analysis, is implemented, tested, and documented in Python. Mature open-source libraries (Ripser, GUDHI, giotto-tda) provide robust infrastructure. Championship broadcast data is secured through the Swansea City AFC--StatsBomb partnership. The primary risk, enhanced GPS data access, is mitigated through multiple independent pathways (Genius Sports UK; Borussia Dortmund, Germany) and the validated broadcast data approach, which demonstrated that framework results are robust without GPS precision. RA recruitment risk is mitigated through early advertising (Month $-2$), competitive salary, and the cross-disciplinary project scope.

\paragraph{Methodology.}
The established multi-scale TDA framework (Brown et al., 2026) operates through six stages: (1)~point cloud construction from 22-player positions; (2)~hierarchical clustering at validated cutoff distances (individual 2.98~m, tactical 12.0~m, team 28.11~m); (3)~adaptive filtration using $\varepsilon_{\max} = \max\bigl(P_{75}(d_{ij}), \max(5.0,\, 2\delta)\bigr)$; (4)~persistent homology computation (Ripser) for $H_0$ and $H_1$; (5)~closed cycle identification through geometric realisation of $H_1$ generators; (6)~statistical analysis (Wilcoxon rank-sum, permutation tests, Spearman correlation, FDR correction). Full methodological detail, including sensitivity analysis and ablation studies, is provided in the ArXiv preprint.

New methods developed in this grant: persistence landscape computation from persistence diagrams at each scale; landscape distance metrics for comparing temporal windows; hierarchical clustering in persistence space for tactical system classification; concurrent processing pipeline for full-season execution.

\paragraph{Previous work and building upon it.}
The ArXiv preprint (Brown et al., 2026) reports comprehensive 11-match validation: three $H_0$ analysis regimes with stability scores ${>}0.88$; 5{,}038 $H_1$ loops detected across 11 matches at two complementary scales; real event correlation (104{,}722 pairs across 10 matches) demonstrating coherent topological response to match dynamics; sensitivity analysis confirming robust $H_1$ detection across a broad cutoff range (8--15~m); and linkage method comparison (single, complete, Ward) characterising chaining effects. Temporal dynamics are match-specific, consistent with the expectation that tactical context drives persistence evolution. This grant builds on the validated framework: from 11 matches to a full season (scaling), from detection to classification (tactical fingerprinting), and from time-series statistics to functional representations (persistence landscapes).

\paragraph{Maximising translation.}
Paper~1 (methodology, Brown et al., 2026) is on ArXiv, with journal submission to JACT during Month~1. Paper~2 (cross-domain application, Brown et al., in preparation) targets submission by Month~4. Paper~3 (full-season results and persistence landscapes) targets Month~11. Conference presentations at SIAM Dynamical Systems and Applied Topology meetings. Open-source code release at Month~12. Industry impact through quarterly meetings with championship club partners. Standard Grant preparation (Months 9--12) incorporates full-season results.

\paragraph{Research environment.}
The project benefits from Swansea University's research environment with established strengths in computational mathematics, data science, and sports analytics. Supercomputing Wales (2~Petaflops) enables ${<}2$\,s per-frame processing, supporting concurrent full-season analysis. Secure GDPR-compliant data storage (100~TB) supports commercial partnerships. The Research Office provides partnership negotiation support, fast-tracked ethics approval (4-week turnaround), and dedicated Impact Officer. Championship broadcast data is secured through the Swansea City AFC--StatsBomb agreement, with GPS data pathways via Genius Sports and Borussia Dortmund. The project is delivered by a Research Associate (1.0~FTE, 9 months, Months 2--10) supervised by PI Dr Rowan Brown, with Co-Investigators Professor Liam Kilduff (Sport and Exercise Sciences) and Professor Gibin Powathil (Mathematics) providing domain expertise and mathematical rigour respectively.
